\documentclass[11pt]{article}
\usepackage[margin=1in]{geometry}
\usepackage{amsmath,amssymb,microtype}
\usepackage[colorlinks=true,urlcolor=blue,citecolor=blue]{hyperref}
\newcommand{\Hh}{\mathcal H}
\newcommand{\D}{\mathcal D}
\newcommand{\Tr}{\operatorname{Tr}}
\title{An affirmative literature answer to the operator $E$-norm question in arXiv:1903.06086}
\author{Literature-status note for \texttt{fa\_banach\_001}}
\date{11 August 2026}

\begin{document}
\maketitle

\noindent\textbf{Status.}
\texttt{literature\_already\_answered}.

\section*{Original question}

Let $G$ be an arbitrary positive densely defined operator on a separable
complex Hilbert space $\Hh$, and let $A$ be $\sqrt G$-bounded on
$\D(\sqrt G)$.  Shirokov's arXiv:1903.06086 asks, immediately after equation
(13) on original PDF page 5, whether the mixed-state operator $E$-norm always
has the simplified expression
\[
 \|A\|_E^G
 =\sup\bigl\{\|A\varphi\|:\varphi\in\D(\sqrt G),\
       \|\varphi\|\leq1,\ \|\sqrt G\varphi\|^2\leq E\bigr\},
\]
without assuming that $G$ is discrete \cite{ShirokovExtension}.

\section*{Exact later answer}

The revised version of Shirokov's \emph{Operator $E$-norms and their use}
answers the question affirmatively \cite{ShirokovENorms}.  Section 5 first
defines $A\rho A^*$ for every finite-energy positive trace-class $\rho$ and
every $\sqrt G$-bounded $A$.  Proposition 6 on original PDF page 20 then
states:
\begin{itemize}
  \item the pure-vector and mixed-state definitions of $\|A\|_E^G$ coincide
        for every $\sqrt G$-bounded $A$; and
  \item both suprema may be taken over subnormalized competitors,
        specifically $\|\varphi\|\leq1$ and $\Tr\rho\leq1$, with energy at
        most $E$.
\end{itemize}
This is exactly equation (13) of the question, with no discreteness assumption
on $G$.  The proof reduces arbitrary finite-energy states to finite-rank
positive operators and uses the concavity of $E\mapsto\|A\|_E^{G\,2}$; the
tail then passes by monotone convergence of $\Tr(A\rho_nA^*)$.

For an independent later cross-check, van Luijk's \emph{Energy-Limited Quantum
Dynamics} gives the subnormalized pure-vector formula as (2.33) and its
equality with the mixed-state supremum as (2.34), in Lemma 2.19(2),(4) on
arXiv PDF pages 15--16 \cite{vanLuijk}.

\section*{Chronology and scope}

The question appeared in arXiv:1903.06086v1 on 14 March 2019.  The answering
arXiv:1806.05668 was revised on 1 October 2020, after that question, and the
published 2020 article contains Proposition 6.  Thus the identifier of the
answering work is numerically earlier, but the decisive revision is later.
No new mathematical result is claimed by this packet.

\begin{thebibliography}{9}
\bibitem{ShirokovExtension}
M.~E. Shirokov,
\emph{On extension of quantum channels and operations to the space of
relatively bounded operators},
Lobachevskii J. Math. \textbf{41} (2020), 714--727;
arXiv:1903.06086,
\url{https://arxiv.org/abs/1903.06086}.

\bibitem{ShirokovENorms}
M.~E. Shirokov,
\emph{Operator $E$-norms and their use},
Sbornik: Mathematics \textbf{211} (2020), 1323--1353;
arXiv:1806.05668v6,
\url{https://arxiv.org/abs/1806.05668}.

\bibitem{vanLuijk}
L.~van Luijk,
\emph{Energy-Limited Quantum Dynamics},
Commun. Math. Phys. \textbf{406} (2025), article 120;
arXiv:2405.10259,
\url{https://arxiv.org/abs/2405.10259}.
\end{thebibliography}

\end{document}
